Le niveau 1 repose sur trois simplifications que le rapport précédent désignait comme ses limites
principales : la chaussée est à surface plane (H6), l'enrobé est homogène à l'échelle du
millimètre (H7) et le contact est sans frottement (H3). Cette section les lève une par une. Le
principe reste le même : chaque modèle ajouté est vérifié contre une solution indépendante, puis
chargé \emph{exactement} par les grandeurs du calcul 3D au point critique, de sorte que son
résultat se raccorde au niveau 1 sans paramètre libre.

\subsection{Frottement pneu--chaussée (levée de H3)}\label{sec:n2_frott}
\paragraph{Modèle.} Au freinage, le pneu exerce sur la chaussée un effort tangentiel dirigé dans
le sens du roulement. On le représente par une loi d'Amontons locale en glissement total :
$q_x(x,y,t)=\mu\,p(x,y,t)$, à l'échelle du pneu \emph{et} à l'échelle de la texture. C'est une
hypothèse majorante pour la texture (en adhérence partielle, le cisaillement local est plus
réparti), et elle est cohérente avec le freinage appuyé, où l'essentiel de l'empreinte glisse. Trois
niveaux : $\mu=0$ (roulement libre, niveau 1), $\mu=0{,}3$ (freinage courant) et $\mu=0{,}6$
(freinage d'urgence sur piste sèche). Par linéarité, la perturbation tangentielle équivalente vaut
$\delta q_{\text{eq}}=\mu\,\delta p_{\text{eq}}$ : la méthode héréditaire (\S\ref{sec:hered}) s'applique
sans modification.

\paragraph{Mise en œuvre.} La réponse nominale avec frottement est calculée par \cs, qui traite
déjà les efforts tangentiels (découplage P-SV / SH). La réponse locale du massif homogène à
$\delta q$ a été ajoutée à \texttt{texture.local\_response}. Pour un massif semi-infini, chaque
amplitude du noyau (normale, tangentielle P-SV, SH) est de la forme $(a+b\,\xi z)\,e^{-\xi z}$
par invariance d'échelle. Les coefficients $(a,b)$ sont identifiés sur le noyau multicouche
général, puis la forme est contrôlée en un troisième point. \textbf{Contrôle T9} : sur un
chargement de texture aléatoire avec $\mu=0{,}5$ et un effort transversal, l'écart à l'évaluation
directe du noyau général vaut \NtNeufHomo{} (massif homogène) et \NtNeufTFE{} (structure complète
du TFE, 5 couches), sur les six composantes et de 1 à \SI{10}{mm}.

\subsection{Modèle local par éléments finis « pixel »}\label{sec:n2_ef}
L'entaille des rainures et la microstructure granulats--mortier sortent du cadre spectral, qui
suppose une couche plane et homogène. Elles sont traitées par un modèle local éléments finis
(\texttt{chausspec/fe2d.py}, écrit pour ce rapport) :
\begin{itemize}[nosep]
\item cellule rectangulaire de largeur $W$, périodique en $x$, de hauteur $H$, encastrée au fond ;
\item maillage régulier de pixels carrés de côté $h$, chacun étant un élément Q4 bilinéaire
      (intégration $2\times2$) portant son propre matériau $(E,\nu)$, $E=0$ désignant un vide ;
\item \emph{déformation plane généralisée} : le déplacement cherché est
      $\bm u=\bm\epsilon^\ast\cdot\bm x+\bm u_c$, où $\bm\epsilon^\ast$ est la déformation nominale
      uniforme ($\epsilon_{xx0}$ dans le plan, $\epsilon_{yy0}$ hors plan) et $\bm u_c$ un champ
      périodique nul au fond. La déformation nominale entre par la précontrainte
      $-\int\bm B^T\bm C:\bm\epsilon^\ast$. Cette formulation « totale » laisse libres les faces
      des vides, ce qu'exige une rainure ;
\item efforts de surface réels (pression et frottement) appliqués sur la face supérieure du
      premier élément solide de chaque colonne ; les fonds de rainure ne sont pas chargés ;
\item $\epsilon_{xy}$ et $\epsilon_{yz}$ nominaux, qui ne sont pas perturbés au premier ordre par
      une géométrie invariante en $y$, sont ajoutés au tenseur avant le calcul de
      $\epsilon_1$ ;
\item résolution directe creuse (factorisation LU unique pour tous les seconds membres).
\end{itemize}

\paragraph{Chargement au point critique.} Pour chaque texture, le calcul 3D (T10) enregistre au
point où $\epsilon_1$ est maximal à \SI{1}{mm} le tenseur nominal complet, la pression
nominale instantanée $p_0$ et la pression nominale héréditaire équivalente
$p_{0,\text{eq}}=E_s\int J(-s)\,\dd_s p_0$. C'est cette dernière qui charge la cellule, modulée par
le contact : $p(x)=p_{0,\text{eq}}\,m(x)$. Dans le cadre (H1)--(H7), cette réduction est exacte
pour la partie texture (\S\ref{sec:hered}). Les déformations nominales $\epsilon_{xx0}$,
$\epsilon_{yy0}$, $\epsilon_{xy0}$ et $\epsilon_{yz0}$ sont prises telles quelles (valeurs
viscoélastiques). \textbf{Conséquence} : la cellule \emph{sans} vide et homogène reproduit le niveau
1. Pour les rainures, qui ne dépendent que de $x$, elle le reproduit exactement : on retrouve
\NplanUn{}~µdef à \SI{1}{mm} contre \NtroisDUn{}~µdef dans le calcul 3D.

\paragraph{Contrôle T11a.} Cellule homogène plane sous la pression de contact des rainures, avec
un frottement $\mu=0{,}5$ et des déformations nominales, comparée à la solution fermée (partie
uniforme analytique, partie périodique par \texttt{local\_response}) : l'écart maximal sur
$\epsilon_{xx}$, $\epsilon_{zz}$ et $\epsilon_{xz}$ vaut \NonzeAgros{} pour $h=\SI{0.2}{mm}$ et
\NonzeAfin{} pour $h=\SI{0.1}{mm}$ à \SI{1}{mm} de profondeur, et moins de
\SI{0.05}{\percent} à \SI{5}{mm}. La convergence est d'ordre 2 environ, comme attendu pour des
éléments Q4. On retient $h=\SI{0.1}{mm}$ (6,35/64).

\paragraph{Distance critique.} Au coin d'une rainure à arêtes vives, la déformation élastique est
singulière : sa valeur au pixel près dépend du maillage et n'a pas de sens physique. On applique
la théorie des distances critiques (Taylor, 2007) : la grandeur pertinente est la déformation
moyennée sur une longueur $L_c$ liée à la microstructure. Pour un enrobé, $L_c$ est de l'ordre de
l'épaisseur des films de mortier et de la taille des plus petits granulats porteurs, soit
\SIrange{0.5}{1}{mm}. Les résultats sont donnés pour $L_c=0$ (pixel), 0,5 et \SI{1}{mm} ; seuls
les deux derniers doivent être interprétés.

\subsection{Entaille des rainures (levée de H6)}\label{sec:n2_rain}
\paragraph{Géométrie.} Un pas de rainurage FAA ($\Lambda=\SI{38.1}{mm}$) avec une rainure
transversale vide de $6{,}35\times\SI{6.35}{mm}$, sur \SI{60}{mm} de profondeur ($W\times H =
384\times605$ pixels de \SI{0.099}{mm}). Deux variantes : arêtes vives, ou arêtes \emph{et} fonds
arrondis de rayon \SI{1}{mm} (arêtes supérieures telles que dans le contact T1, congés de fond
obtenus en remplissant les coins par un quart de cercle). La coupe est dans le plan $(x,z)$ du
roulement, perpendiculaire aux rainures : c'est le plan où la rainure interrompt la couche.
\paragraph{Chargement.} Pression de contact des plateaux (T1, $\bar p=\SI{1.65}{MPa}$) mise à
l'échelle de $p_{0,\text{eq}}$ au point critique « rainures » de T10, frottement éventuel
$\mu p$, déformations nominales du même point. La référence est la même cellule sans vide
(niveau 1). On note $K_{H6}(z)$ le rapport des maxima de $\epsilon_1$ à la profondeur $z$ sous le
plateau.

\subsection{Granulats et mortier (levée de H7)}\label{sec:n2_mortier}
\paragraph{Microstructure.} Coupe 2D de la couche de roulement (\SI{40}{mm} $\times$ \SI{40}{mm},
pixels de \SI{0.1}{mm}) : granulats circulaires de 2 à \SI{10}{mm}, tirés classe par classe du plus
gros au plus petit selon une courbe de Fuller (passant $\propto(d/D)^{0{,}45}$), placés par
tirage séquentiel aléatoire avec un jeu minimal de \SI{0.2}{mm} entre granulats et une
périodicité en $x$. La fraction surfacique obtenue est de 0,53 à 0,56, ce qui est typique des
coupes d'enrobés 0/10 pour la fraction supérieure à \SI{2}{mm}. Le reste est le \emph{mortier
bitumineux} (liant, fines et sables de moins de \SI{2}{mm}), traité comme un milieu homogène. La
surface est sciée : les granulats y sont tronqués. Trois tirages indépendants par cas.
\paragraph{Homogénéisation et mise à l'échelle.} Une cellule biperiodique de même granulométrie,
soumise à trois déformations uniformes, donne le tenseur effectif $\bm C_{\text{hom}}$ dont on
extrait $(E_{\text{hom}},\nu_{\text{hom}})$ isotropes. Les modules des phases sont ensuite mis à
l'échelle pour que $E_{\text{hom}}$ soit le module de l'enrobé équivalent du niveau 1. Seul compte
alors le \emph{contraste} $c=E_{\text{granulat}}/E_{\text{mortier}}$. Avec $\nu_g=0{,}25$ et
$\nu_m=0{,}41$, on obtient $\nu_{\text{hom}}\approx0{,}35$ comme au niveau 1. Le contraste dépend
de la température et du temps de charge : de l'ordre de 10 pour un mortier froid et rapide,
supérieur à 100 pour un mortier chaud sous charge lente (granulats de 50 à \SI{70}{GPa} contre un
mortier de 0,5 à \SI{5}{GPa}). On étudie $c=10$, 30 et 100.
\paragraph{Chargement et indicateurs.} Point critique de la texture aléatoire MPD \SI{1.0}{mm}
(T10) : ligne de pression de contact extraite de la carte 3D à $\bar p=\SI{1.65}{MPa}$ et mise à
l'échelle de $p_{0,\text{eq}}$, plus les déformations nominales. Une ligne de taches de contact
prolongée indéfiniment en $y$ charge plus qu'en 3D : l'amplitude de la modulation est étalonnée
pour que la cellule homogène reproduise $\epsilon_1$ du calcul 3D à \SI{1}{mm} (voir les
résultats ; la ligne brute sert d'étude de sensibilité). On calcule aussi la même cellule
sous une pression \emph{uniforme} $p_{0,\text{eq}}$ (surface lisse, TFE) et sous les seules
déformations nominales. La référence est la cellule homogène $(E_{\text{hom}},\nu_{\text{hom}})$
sous le même chargement. Indicateur : facteur de concentration dans le mortier
$K_m(z)=\max\epsilon_1^{\text{mortier}}/\max\epsilon_1^{\text{homogène}}$ dans une bande de
$\pm\SI{0.25}{mm}$ autour de $z$ (maximum et quantile à 99\,\%).
